% Generated draft: requires mathematical and source review.
\documentclass[11pt]{article}
\usepackage[margin=25mm]{geometry}
\usepackage{fontspec}
\setmainfont{Latin Modern Roman}
\usepackage{amsmath,amssymb,mathtools,unicode-math}
\setmathfont{Latin Modern Math}
\usepackage{enumitem,needspace,xurl,hyperref}
\setlength{\parindent}{0pt}
\setlength{\parskip}{5pt}
\setlength{\emergencystretch}{3em}
\newcommand{\N}{\mathbb{N}}
\newcommand{\Z}{\mathbb{Z}}
\newcommand{\Q}{\mathbb{Q}}
\newcommand{\R}{\mathbb{R}}
\newcommand{\C}{\mathbb{C}}
\newcommand{\F}{\mathbb{F}}
\newcommand{\A}{\mathbb{A}}
\newcommand{\PP}{\mathbb{P}}
\newcommand{\E}{\mathbb{E}}
\newcommand{\eps}{\varepsilon}
\newcommand{\dd}{\,\mathrm{d}}
\begin{document}
\section*{UnsolvedMath ID: 30004492}
\textbf{Source problem code:} catalogue entry
\subsection{Definitions and mathematical statement}
Let \(B\) be a smooth projective variety, let \(Z\subset B\) be a reduced normal-crossing divisor, and set \(U=B\setminus Z\). Here reduced means that each component of the divisor has multiplicity one, and the normal-crossing condition means that locally its branches are coordinate hyperplanes. Suppose that a period map
\[
\Phi:U\longrightarrow\Gamma\backslash D
\]
is given, and write \(\mathcal P=\Phi(U)\). The notation \(D\) denotes the period domain in the supplied statement, and \(\Gamma\) denotes the group occurring in its quotient; their precise defining data are not supplied.

\textbf{Problem.} For these data, determine what natural completions of \(\Phi\) to maps defined on all of \(B\) exist, and determine their properties. Schematically, an extension to a target \(\overline{\mathcal P}\) containing \(\mathcal P\), with inclusion \(j:\mathcal P\hookrightarrow\overline{\mathcal P}\), would satisfy
\[
\overline\Phi:B\longrightarrow\overline{\mathcal P},
\qquad
\overline\Phi|_U=j\circ\Phi_{\mathcal P},
\]
where \(\Phi_{\mathcal P}:U\to\mathcal P\) is \(\Phi\) with codomain restricted to its image. Thus every point of \(Z\) must lie in the domain of the sought extension.

This diagram records only the extension requirement, not a complete definition of the requested completion. The source does not specify the category of the target and maps, the topology or geometric structure on \(\mathcal P\), the density requirement, or a criterion making a completion natural. It also does not state the base field or the detailed hypotheses on the period map. Accordingly, the requested classification and properties remain qualified by these unresolved choices; no particular algebraic or topological completion theorem is asserted here.
\subsection{Short English statement}
Given a period map on the complement of a reduced normal-crossing divisor in a smooth projective variety, determine natural ways to extend it over the entire variety by completing its image, and describe the resulting properties. The intended category, period-map hypotheses, and meaning of natural completion require clarification.
\subsection{Sources}
\begin{itemize}
\item Local UnsolvedMath registry, supplied complete source record: permanent ID \(30004492\), title \emph{Natural Completions of Period Maps}, exact source problem code \texttt{catalogue entry}, in the Oberwolfach Reports Open Problems collection. The record gives proposed year 2020 and names no individual proposer.
\item Supplied report citation: \emph{Komplexe Analysis -- Algebraicity and Transcendence} (2021), Oberwolfach Reports, pp.~1263--1313, DOI \texttt{10.4171/owr/2020/24}. \url{https://doi.org/10.4171/owr/2020/24}. The supplied literature-source label is \emph{Oberwolfach Report 24/2020}. The differing year indications are retained without reconciliation.
\item Supplied literature source: \emph{Global asymptotic structure of period mappings}. \url{https://arxiv.org/abs/2010.06720}. No authors or specific theorem references are supplied in the record.
\item Supplied dated literature curation, checked \texttt{2026-08-22}: the record labels the problem partially solved and reports that a proper topological SBB-type completion with projective fibers and quasi-projective image strata has been constructed, while the general algebraic completion and its full properties remain incomplete. This is the supplied dated assessment derived from OWR curation, not an independently verified mathematical or current-literature assessment. Its time-sensitive status and cited claims require independent verification.
\end{itemize}
\end{document}
